\appendixsection{Ключевые промпты узлов графа и схема их сборки}

Модуль \texttt{prompts.py} содержит порядка двадцати задачных промптов~--- по одному (а иногда нескольким) на каждый содержательный узел графа. Все промпты следуют одному шаблону: системный пролог \texttt{LEGAL\_ASSISTANT\_SYSTEM\_PROMPT}, задаваемая узлу инструкция с явно перечисленными правилами и (для большинства узлов) точная спецификация JSON-ответа, который модель должна вернуть в режиме структурированного вывода. Типовой пример промпта и схема его сборки в узле графа приведены ниже.

Системный пролог \texttt{LEGAL\_ASSISTANT\_SYSTEM\_PROMPT} (рисунок~\ref{src:prompt-system}) подключается ко всем содержательным вызовам языковой модели и задаёт общую роль <<юрист, специализирующийся на российском гражданском праве>>, формат упоминания источников и запрет на выдуманные факты.

\begin{figure}[H]
\begin{lstlisting}[style=promptstyle]
Ты -- юрист, специализирующийся на российском гражданском праве.
Тебя зовут PactumAI.
Источники информации в каждом запросе:
- описание сделки и история диалога -- это то, что сообщил пользователь;
- разделы "Общие нормы", "Специальные нормы", "Нормы" -- это выдержки из
  законодательства РФ, автоматически найденные в базе нормативных актов;
  пользователь их не присылал и не видит.
Не используй формулировки "в предоставленных материалах", "вы предоставили
нормы", "вы дали ссылки" -- нормы взяты из законодательства РФ. Используй
формулировки "согласно законодательству РФ", "в соответствии с [источник,
статья]".
Отвечай только на основе предоставленных норм и описания сделки.
Не выдумывай факты.
\end{lstlisting}
\caption{Системный промпт LEGAL\_ASSISTANT\_SYSTEM\_PROMPT}
\label{src:prompt-system}
\end{figure}

В качестве содержательного примера приведён промпт генерации проекта договора \texttt{GENERATE\_CONTRACT\_PROMPT} (рисунок~\ref{src:prompt-generate})~--- наиболее ответственного узла графа, формирующего основной артефакт системы. Промпт включает HTML-шаблон делового документа с фиксированной структурой из девяти разделов, явно перечисляет разрешённые HTML-теги (\texttt{body}, \texttt{p}, \texttt{strong}, \texttt{table}, \texttt{tr}, \texttt{td}), запрещает использование Markdown и пояснений вне разметки и фиксирует ключевое правило <<в самом тексте договора не упоминать статьи закона>>: формулировки условий пишутся обычным языком, а соответствующие нормы и риски выносятся в раздел рекомендаций отдельным узлом.

\begin{figure}[H]
\begin{lstlisting}[style=promptstyle]
Сгенерируй только HTML-тело договора: от <body> до </body>. Договор должен
соответствовать стандартному шаблону по оформлению (9 разделов: предмет, срок,
права и обязанности, цена, ответственность, изменение/расторжение, порядок
разрешения споров, заключительные положения, реквизиты и подписи).

Обязательно добавь в договор поля для каждой из сторон: ФИО (физлицо),
ОГРНИП/ИНН (ИП), полное наименование/ОГРН/ИНН (юрлицо), адрес.

Правила:
- Не выдумывай факты и реквизиты.
- Для недостающих данных используй только печатные поля: "______", "20__".
- Разрешены только теги: body, p, strong, table, tr, td.
- Без markdown/json/пояснений.
- Не указывай в тексте договора ссылки на статьи закона, кодексы или иные
  нормативные акты. Формулировки условий пиши обычным языком.
\end{lstlisting}
\caption{Промпт генерации проекта договора GENERATE\_CONTRACT\_PROMPT}
\label{src:prompt-generate}
\end{figure}

Схема сборки промпта в узле графа единообразна и иллюстрируется на примере узла \texttt{check\_written\_form} (рисунок~\ref{src:prompt-assembly}). Системный пролог \texttt{LEGAL\_ASSISTANT\_SYSTEM\_PROMPT} и инструкция узла \texttt{CHECK\_WRITTEN\_FORM\_PROMPT} являются константами; в момент вызова к ним конкатенацией добавляется человеческий промпт, содержащий релевантные поля состояния агента~--- историю диалога, тип сделки, исходное описание, найденные нормы. Названия секций (<<История диалога:>>, <<Тип сделки:>>, <<Описание сделки (от пользователя):>>, <<Нормы (выдержки из законодательства РФ):>>) вынесены в отдельные константы, что обеспечивает единообразие между промптами и согласованность с инструкцией системного пролога о способах ссылок на содержимое разделов. Вызов \texttt{\_invoke\_json} обращается к LLM в режиме структурированного вывода с предопределённой fallback-схемой на случай ошибки парсинга и возвращает уже разобранный объект, готовый к использованию в логике маршрутизации.

\begin{figure}[H]
\begin{lstlisting}[style=promptstyle, language={}]
def check_written_form(state: ContractAgentState) -> dict:
    fallback = {"requires_written_form": True,
                "reason": "письменная форма предполагается по умолчанию"}
    result = _invoke_json(
        "check_written_form",
        CHECK_WRITTEN_FORM_PROMPT,                      # системный + инструкция узла
        (                                                # человеческий промпт:
            f"{SECTION_LABEL_HISTORY}\n"                 #   - история диалога
            f"{_conversation_context(state)}\n\n"
            f"{SECTION_LABEL_DEAL_TYPE} "                #   - тип сделки
            f"{state.get('deal_type') or '-'}\n\n"
            f"{SECTION_LABEL_DEAL_DESCRIPTION}\n"        #   - описание от пользователя
            f"{state.get('deal_description', '')}\n\n"
            f"{SECTION_LABEL_NORMS}\n"                   #   - найденные нормы
            f"{_norms_to_context(state.get('retrieved_norms'))}"
        ),
        fallback=fallback,
    )
    ...
\end{lstlisting}
\caption{Схема сборки промпта в узле графа (фрагмент функции \texttt{check\_written\_form} из модуля \texttt{nodes.py})}
\label{src:prompt-assembly}
\end{figure}

Остальные содержательные промпты модуля \texttt{prompts.py} построены по той же схеме и решают аналогичные задачи в своих узлах графа: \texttt{build\_classify\_deal\_prompt} (классификация типа сделки), \texttt{CHECK\_GENERAL\_NORMS\_PROMPT} (проверка препятствий по общим нормам), \texttt{BUILD\_RETRIEVAL\_QUERY\_PROMPT} (формирование поискового запроса для подсистемы RAG), \texttt{FILTER\_RETRIEVED\_NORMS\_PROMPT} (LLM-фильтр выдачи), \texttt{CHECK\_WRITTEN\_FORM\_PROMPT} (необходимость письменной формы), \texttt{build\_check\_data\_sufficiency\_prompt} (достаточность данных), \texttt{VALIDATE\_CONTRACT\_PROMPT} (валидация проекта договора), \texttt{GENERATE\_RECOMMENDATIONS\_PROMPT} (формирование юридического отчёта), \texttt{DETECT\_VAGUE\_REFERENCES\_IN\_NORMS\_PROMPT} и \texttt{INTEGRATE\_ENRICHMENTS\_PROMPT} (обработка бланкетных отсылок), \texttt{EDIT\_CONTRACT\_PROMPT} (точечная правка договора), \texttt{CLASSIFY\_FOLLOWUP\_INTENT\_PROMPT} и \texttt{FOLLOWUP\_TRIAGE\_PROMPT} (маршрутизация дополнительных вопросов), \texttt{FOLLOWUP\_SUBAGENT\_PROMPT} (изолированный поиск субагента), \texttt{CLASSIFY\_YESNO\_PROMPT} (классификация ответа <<да/нет>>), \texttt{build\_inform\_unsupported\_deal\_prompt} и \texttt{INFORM\_USER\_PROMPT} (формирование вежливых сообщений об отказе), \texttt{GENERATE\_CASE\_TITLE\_PROMPT} (формирование заголовка кейса для интерфейса). Тексты промптов хранятся в модуле \texttt{app/agent/prompts.py} и снабжены подробными встроенными комментариями, поясняющими каждое правило.
